\section{Задача и модель}
\begin{definition}[Сценарий обработки]
Реализована консольная модель на C11/POSIX в одном процессе и потоке. Посылка проходит входную линию, накопитель, сканер, сортировщик, конвейер и выход нужного направления. Накопители ограничены; заполнение следующей очереди блокирует устройство без потери посылки. После исчерпания повторов чтения посылка направляется на ручную обработку.
\end{definition}

Параллелизм имитируется циклом по тактам: разгрузка, конвейеры, сортировщики, сканеры, источники. Устройства хранят посылку и время готовности; поступления и ошибки случайны. На направление выделен один конвейер; время перевозки равно $\lceil L/v\rceil$. Выходы разгружаются периодически; доставка учитывается при поступлении в выход. Ручная обработка не моделируется.

Структура Parcel хранит номер, направление, число попыток и стадию посылки. Устройство обслуживает одну посылку; очереди соблюдают FIFO. Входная очередь общая, перед конвейерами очереди отдельные. Порядок обновления устройств не позволяет посылке пройти несколько этапов за такт.

\section{Интерфейс и запуск}
Аргументы задают объём партии, устройства, вместимости, длительности, ошибки, повторы и seed; полный список доступен через \texttt{-{}-help}. Без аргументов обрабатываются 20 посылок. Консоль показывает такт, номер посылки, событие и итоговую статистику. Журнал дублирует вывод через \texttt{open/write/close}.

Параметр \texttt{-{}-seed} воспроизводит события в той же среде; \texttt{-{}-delay-ms} меняет только скорость показа. Журнал перезаписывается; запись учитывает частичный вывод и EINTR.
\begin{example}[Команды запуска]
\begin{lstlisting}[basicstyle=\small\ttfamily,frame=none,backgroundcolor=\color{example_back},aboveskip=0pt,belowskip=0pt,xleftmargin=0pt,xrightmargin=0pt]
make
./build/sorting_center --parcels 10 --error-percent 100 \
    --retries 2 --log run.log
make test
\end{lstlisting}
\end{example}
\section{Завершение работы}
\begin{remark}[Остановка и сохранение состояния]
После доставки или передачи на ручную обработку всей партии возвращается код 0. Ctrl+C/SIGTERM дают коды 130/143 и частичную статистику; память освобождается, журнал закрывается. Сохраняется баланс: запланировано = доставлено + отказы + незавершённые + ещё не созданные. Некорректные аргументы отклоняются с кодом 2.

Обработчик сигнала только устанавливает флаг. Основной цикл выводит итоги и освобождает ресурсы; сигнал также прерывает ожидание в nanosleep. Ошибки вывода возвращают код 1.

\end{remark}

\clearpage
\section{Проверка программы}
\begin{solution}[Результаты тестирования]
На macOS, Apple Clang 17 прошли 14 групп тестов в обычном режиме и с \texttt{-{}-check}. Из 10 посылок при 0\% ошибок доставлены все; при 100\% и двух повторах все переданы вручную за 30 чтений. При вместимости 1 обработаны все 20 посылок; пустая партия завершается сразу. Оба сигнала прерывают длительную задержку с сохранением итогов.

По журналу тесты восстанавливают маршрут каждой посылки и проверяют вместимости, направление и длительности. Проверены 35 наборов параметров и переходы кольцевой очереди через конец массива при вместимостях 1, 3 и 7.

\end{solution}

\section{Исправления и оценка}
\begin{exercise}[Исправления и ограничения]
По вопросам из \path{docs/DEBUGGING.md}: два повтора означают три попытки; выходы освобождаются без повторного учёта доставки; очередь сделана кольцевой с операциями $O(1)$, обход учитывает head. Полный контроль каждого такта включается через \texttt{-{}-check}, перед завершением выполняется всегда. Сложность основного цикла $O(TK+N)$, с проверкой каждого такта $O(T(N+K))$, где $T$ - такты, $K$ - источники и устройства, $N$ - посылки; паузы и ожидание вывода не учтены. Недостаток - перебор пустых тактов.

Обход очереди использует индекс \texttt{(head + i) \% limit}, иначе после смещения начала проверка читала бы чужие элементы.

\end{exercise}

\section{Использование ИИ}
Четыре запроса из \path{docs/PROMPTS.md}:
\begin{enumerate}[leftmargin=*,itemsep=0.6em,topsep=0.4em]
\item У меня retries = 2, и после второй ошибки посылка уходит на ручную обработку. Правильно ли я понимаю, что нужны первое чтение и ещё два повтора? Подскажи, с чем сравнивать attempts перед повтором
\item После удаления первого элемента я сдвигаю всю очередь. Можно ли обойтись без сдвига, оставив обычный массив? Объясни на маленьком примере
\item Переделал очередь в кольцевую. Есть head и count, но проверка всё ещё обходит items[i] от нуля до count. После нескольких извлечений она сообщает об ошибке. Что я не учёл в обходе?
\item После каждого такта я обхожу все созданные посылки. Верно ли, что при T тактах и N посылках получается O(T * N)? Если полный обход оставить под флагом \texttt{-{}-check}, какие проверки нужны в обычном режиме?
\end{enumerate}
